\documentclass{MITcsail}

\title{Artificial Intelligence in Alzheimer's Disease Research}
\author{AD Literature Knowledge Pipeline\\
\vspace{1em}
\normalfont{\small Automated narrative literature review}}

\begin{document}

\maketitle
\thispagestyle{firstpagestyle}

\begin{abstract}
The application of artificial intelligence (AI) in Alzheimer's disease (AD) research is gaining prominence, addressing critical challenges in diagnosis, progression monitoring, and treatment personalization. This narrative literature review synthesizes findings from 83 studies published between 2020 and 2025, highlighting the significant role of deep learning and machine learning techniques in improving detection and prediction outcomes. Key findings indicate that various AI models achieved accuracies up to 99.05\% for classifying dementia stages, showcasing their potential for early AD diagnosis. However, considerable methodological inconsistencies and a lack of standardized frameworks hinder reproducibility across studies. Moreover, gaps in dataset diversity and insufficient exploration of multimodal integrations are evident. This review emphasizes the need for advancing AI methodologies in AD research, suggesting that future work should prioritize reproducibility and validation across diverse populations to optimize diagnostic precision.
\end{abstract}

\section{Introduction}

Alzheimer's disease (AD), a progressive neurodegenerative disorder, affects millions worldwide and leads to significant cognitive decline and loss of functional independence. The escalating prevalence of AD, alongside its profound implications for individuals, families, and healthcare systems, underscores the urgent need for early detection and effective management strategies. As the artificial intelligence (AI) domain advances, it presents innovative opportunities for enhancing our understanding of and response to AD. By integrating AI methodologies, particularly in diagnostics and biomarker discovery, substantial opportunities arise to advance research and clinical outcomes \citep{ref_10_1016_j_media_2020_101694,ref_10_1002_alz_13390}. This review aims to synthesize and evaluate the literature centered on AI applications in Alzheimer's disease, exploring methodologies such as deep learning, machine learning, and neuroimaging analysis. The papers considered were published between January 2020 and June 2025, ensuring a contemporary outlook on innovations and challenges in this area. Key topics encompass diverse methodologies, the evolution of diagnostic frameworks, and multi-modal data integration \citep{ref_10_3389_fpubh_2022_853294,ref_10_3389_fpubh_2022_834032}. The review is structured to outline the broader implications of AI in AD research, followed by in-depth analyses of primary methodological approaches and their applications. The discussion concludes with limitations and future research directions. Despite the expanding body of work, gaps persist, particularly regarding reproducibility across studies and the necessity for diverse datasets to enhance the generalizability of AI models \citep{ref_10_1016_j_media_2020_101694,ref_10_1016_j_cmpb_2021_106032}. This review underscores the potential of AI-driven technologies to transform Alzheimer's disease research and management.

\section{Review Methodology}

\textit{A narrative review synthesizing literature on artificial intelligence in Alzheimer's disease, identifying methodological patterns and gaps in current research.}

This literature review was structured according to a narrative methodology, encompassing 83 papers published from January 2020 to June 2026, all examining various applications of artificial intelligence (AI) in Alzheimer's disease research. The selection yielded 83 usable papers with none excluded due to quality, ensuring a robust evidence base for synthesis. The review focused on comprehensively identifying methodologies used, study designs employed, reported findings, and articulated limitations in the selected literature. Key methodologies included deep learning—particularly convolutional neural networks—machine learning, and ensemble methods, with model development studies constituting the majority of papers. Synthesized findings pointed to challenges such as variability in participant selection, data privacy concerns, and reproducibility issues due to a lack of standardized frameworks. Despite extensive exploration, significant gaps remain, particularly regarding the generalizability of AI models across diverse populations and the need for comprehensive datasets that reflect broad cohorts. These findings highlight the continuing need for innovative approaches in both research design and the application of AI technologies in clinical settings to enhance early diagnosis and treatment strategies for Alzheimer's disease.

\section{Background and Related Literature}

\subsection{Alzheimer's Disease In The Included Literature}

\textit{This section discusses the multifaceted applications of AI in Alzheimer's disease research, synthesizing findings across various studies that employ advanced computational methods such as deep learning and machine learning for diagnosis and prognosis. The literature underscores the significance of early detection, which correlates with more effective treatment outcomes. Key methodologies, findings, and identified research gaps are highlighted, providing a coherent overview of the present state and future needs in this domain.}

Artificial Intelligence (AI) has emerged as a critical tool in understanding, diagnosing, and managing Alzheimer's disease (AD), a progressive neurodegenerative disorder that impairs memory and cognitive functions. The urgency of developing effective diagnostic and therapeutic strategies in Alzheimer's research is amplified by the disease's rising prevalence globally. Notably, early intervention is posited to improve treatment outcomes, making early diagnosis paramount \citep{ref_10_3389_fpubh_2022_853294}. AI-based methodologies, ranging from deep learning to machine learning models, have been pivotal in enhancing diagnostic accuracy and prognostic capabilities. The literature reveals diverse applications of AI frameworks in AD, with a significant emphasis on deep learning models. For instance, convolutional neural networks (CNNs) have demonstrated exemplary performance in identifying and classifying stages of Alzheimer's disease from neuroimaging data such as MRI (Nanni et al., 2020; Murugan et al., 2021). Multiple studies report high accuracy rates, such as up to 96\% in CAD systems, validating the effectiveness of AI in clinical settings \citep{ref_10_3389_fpubh_2022_834032}. Despite these promising results, methodological inconsistencies and generalizability challenges underscore the need for standardized frameworks in AI studies \citep{ref_10_1016_j_media_2020_101694}. Moreover, the diversity and representativeness of datasets utilized in AI training models pose significant barriers, with ongoing calls for larger and more inclusive datasets to refine model performance \citep{ref_10_1002_alz_13390}. Additionally, issues related to reproducibility, where some studies lack publicly accessible implementation details, highlight critical gaps in the literature. Addressing these gaps could significantly enhance the reliable integration of AI technologies in everyday clinical practice for Alzheimer's diagnosis and management. In summary, while AI presents transformative promise in Alzheimer's disease research, ongoing efforts are necessary to improve methodological rigor, dataset diversity, and reproducibility of results to optimize outcomes across varied patient populations.

\subsection{Artificial Intelligence In The Included Literature}

\textit{This section discusses the application and implications of artificial intelligence (AI) methods in Alzheimer's research, emphasizing their impact on diagnosis, prediction, and treatment.}

Artificial intelligence encompasses a diverse array of methods that significantly enhance the capabilities of Alzheimer's disease (AD) research. Within the literature, deep learning techniques, particularly convolutional neural networks (CNNs), are prominently featured for imaging analysis in AD. For instance, models like DEMNET achieve high accuracy (95.23\%) in diagnosing early-stage AD using MR images \citep{ref_10_1109_access_2021_3090474}. Furthermore, AI-driven approaches for biomarker discovery may expedite the identification of dementia-related biomarkers \citep{ref_10_1002_alz_13390}. However, significant gaps in standardization and reproducibility of AI methodologies remain. Many existing frameworks lack accessibility and detailed implementation guidelines, complicating comparative analyses \citep{ref_10_1016_j_media_2020_101694}. This underscores the ongoing need for standardized protocols within the research community to enhance model reproducibility and address discrepancies in classification outcomes. Overall, while AI methods exhibit promising results in diagnosing and stratifying Alzheimer's disease, limitations such as varied methodologies and insufficient generalizability across distinct populations must be acknowledged, directing future research towards addressing these challenges.

\section{Methods and Analytical Approaches}

\subsection{Methodological Landscape}

\textit{This section delineates the methodology landscape in artificial intelligence applications for Alzheimer's disease research, categorized by various computational techniques and associated study designs.}

The methodological landscape within the application of artificial intelligence (AI) in Alzheimer's disease (AD) research encompasses a diverse array of computational approaches. The literature reveals a substantial emphasis on deep learning techniques, specifically convolutional neural networks (CNNs) and more generalized deep learning methods, with 33 and 29 papers, respectively, dedicated to these methodologies. This focus indicates a growing reliance on advanced modeling strategies to enhance diagnostic accuracy. Moreover, traditional machine learning methods, including support vector machines (SVMs) and random forests, have also been employed, with 19 and 1 papers respectively. Conversely, ensemble learning has played a notable role in 8 studies, demonstrating a commitment to composite model architectures aimed at improving diagnostic outcomes. The methodological variety extends to emerging techniques such as transfer learning and graph neural networks, significant even though they are represented by only 4 and 2 papers respectively. This trend indicates an innovative shift towards leveraging pre-trained models to adapt to the specific nuances of AD diagnostics. Diverse study designs are reflected in the literature, with a predominant focus on model development studies (37 papers), alongside validation studies (16 papers) that confirm the effectiveness of these AI models across various cohorts. Observational studies (6 papers) and cohort studies (4 papers) further illustrate real-world applications in Alzheimer's diagnostics and monitoring. Nonetheless, several papers remain categorized as exhibiting unclear study designs (10 papers), highlighting the necessity for improved methodological clarity in future research. Overall, the landscape of AI methodologies in Alzheimer's research is rich and complex, emphasizing the importance of robust validation to support clinical utility and adoption.

\subsection{Comparison of Approaches}

\subsection{Evidence Patterns Across Approaches}

\textit{This section synthesizes the convergences, divergences, and uncertainties in findings across various methodological approaches employed in Alzheimer's disease research using artificial intelligence. A comprehensive evaluation reveals that while there is significant progress in using deep learning and ensemble techniques for diagnosing Alzheimer's, variations in methodology contribute to inconsistencies in reporting overall performance metrics.}

The deployment of artificial intelligence (AI) in Alzheimer's disease research has yielded diverse findings across different methodological approaches. Deep learning models, notably convolutional neural networks (CNNs), have been extensively utilized for classification and prediction of Alzheimer's disease progression. Studies by \citet{ref_10_1016_j_media_2020_101694} and \citet{ref_10_3389_fpubh_2022_834032} highlight a critical issue: the difficulty in comparing classification performances stemming from methodological inconsistencies and potential data leakage. This concern accentuates a recurring theme in the literature where reproducibility is hampered by a lack of accessible frameworks and detailed implementation protocols \citep{ref_10_1016_j_media_2020_101694}. The accuracy of models deployed through deep learning and machine learning techniques varies considerably. For instance, \citet{ref_10_1109_access_2021_3090474} achieved a notable accuracy of 95.23\% with the DEMNET model, while others reported a wide range of success typical for this field. Proposed models have effectively attained validation accuracies from 83\% to 99\% in distinguishing Alzheimer's patients from healthy controls, often leveraging expansive datasets and multiple imaging modalities \citep{ref_10_3389_fpubh_2022_853294,ref_10_1016_j_mri_2021_02_001}. However, limitations exist, as highlighted by various authors, where smaller sample sizes or insufficient diversity in datasets restrict the generalizability of these AI models. Despite the promising results observed from AI methodologies, multiple studies emphasize the need to integrate varied types of clinical data to bolster performance and interpretability of these models \citep{ref_10_1093_jamia_ocac168,ref_10_1016_s2589_7500_22_00169_8}. As consistently noted, future work must focus on enhancing data integration and dataset diversity to fortify the roles of these AI approaches in the early diagnosis and ongoing monitoring of Alzheimer's disease.

\section{Data Foundations and Study Designs}

\subsection{Datasets and Data Sources}

\textit{This section examines the datasets and data sources utilized in the literature concerning AI applications in Alzheimer's disease research. A variety of data origins are analyzed, ranging from standard publicly available datasets to newly collected data and proprietary sources.}

In recent studies, various datasets have been employed to explore the effectiveness of artificial intelligence methods in diagnosing and monitoring Alzheimer's Disease (AD). A predominant data source derives from neuroimaging modalities such as Magnetic Resonance Imaging (MRI) and Positron Emission Tomography (PET). For instance, \citet{ref_10_3389_fpubh_2022_834032} utilized 18FDG-PET images involving 855 patients, underscoring the significance of robust datasets in evaluating AI models' performance. Similarly, MRI data significantly contributes to research efforts, impacting model training and validation \citep{ref_10_1016_j_neuroimage_2020_117203,ref_10_1109_access_2021_3090474}. Moreover, datasets vary in their accessibility and diversity. The Alzheimer's Disease Neuroimaging Initiative (ADNI) has provided crucial resources, enhancing the reproducibility of research outputs. Limitations remain in achieving consistency across studies, as highlighted by \citet{ref_10_1016_j_media_2020_101694}, who emphasized challenges related to data leakage and methodological discrepancies affecting classification performance. Furthermore, the lack of standardized frameworks significantly hinders comparability of results among different studies \citep{ref_10_1016_j_media_2020_101694,ref_10_1002_alz_13390}. Several studies call for further investigation into improving data integration methods and utilizing additional biomarkers for enhanced diagnostic accuracy \citep{ref_10_1002_alz_13390,ref_10_1016_s2589_7500_22_00169_8,ref_10_3389_fpubh_2022_772592}. This highlights the need for concerted efforts to establish comprehensive and standardized datasets that can support future advancements in Alzheimer's research and AI applications. Overall, the datasets utilized range from public repositories like ADNI to proprietary datasets, with many studies indicating the necessity for diversified sources to enhance model generalizability and diagnostic accuracy.

\subsection{Samples, Cohorts, and Study Populations}

\textit{This section synthesizes the various samples, cohorts, and populations studied in papers focusing on AI applications in Alzheimer's disease research. The distribution of samples across different methodologies and designs highlights important patterns in data representation and cohort characteristics, reflecting the diversity and limitations inherent in current research designs.}

Recent studies have utilized diverse samples and cohorts to investigate the application of artificial intelligence (AI) in Alzheimer's disease research. Sample sizes vary significantly, with some studies involving hundreds of participants, while others focus on smaller cohorts, affecting the generalizability of findings \citep{ref_10_1016_j_media_2020_101694,ref_10_3389_fpubh_2022_853294}. A noteworthy study assessing deep learning models used 18FDG-PET images from 855 patients \citep{ref_10_3389_fpubh_2022_834032}. Emerging methodologies utilize different imaging techniques, such as MRI and PET scans, to analyze brain structure and function in Alzheimer's patients. This approach underscores the potential for integrating multimodal imaging data to enhance diagnostic accuracy. Many research designs target early detection and prediction, relying on longitudinal data from participants at various stages of the disease, which informs tailored therapeutic approaches \citep{ref_10_3389_fpubh_2022_853294,ref_10_3389_fpubh_2022_834032}. However, limitations exist within current cohort studies, including biases related to sample selection and the settings in which the data is collected. Research often focuses on homogeneous groups, raising concerns about the applicability of results to broader and more diverse populations, highlighting the need for greater representation in future studies \citep{ref_10_1002_alz_13390,ref_10_3389_fpubh_2022_772592}. Additionally, AI models' generalizability is often constrained by the specific characteristics of training datasets, necessitating further validation across varied demographic groups to ensure robustness in clinical settings \citep{ref_10_1016_j_neuroimage_2020_117203}.

\subsection{Study Designs and Validation Strategies}

\textit{This section compares various study designs and validation strategies across the literature on artificial intelligence in Alzheimer's disease research. It highlights the dominant use of model development studies and varying methodologies, elucidating how these factors shape evidential strength and comparability across studies. Notably, methodological inconsistencies and potential biases are observed, marking significant gaps in reproducibility and external validation.}

In the realm of artificial intelligence applied to Alzheimer's disease, the studies reviewed exhibit diverse designs and validation strategies. Among the 83 papers considered, 37 focus on model development, while 16 emphasize validation studies. Observational and cohort studies comprise six and four respective contributions, underscoring a prevalent inclination towards creating new AI models for diagnosis and prediction \citep{ref_10_1016_j_media_2020_101694,ref_10_1016_j_neuroimage_2020_117203}. The methodologies employed primarily include deep learning and convolutional neural networks, accounting for 33 and 29 papers, respectively. Other methods such as machine learning, ensemble learning, and transfer learning also feature prominently, albeit to a lesser extent. This strong leaning towards deep and convolutional methods reflects current trends in neural network applications \citep{ref_10_3389_fpubh_2022_834032,ref_10_3389_fpubh_2022_853294}. However, the variation in study designs presents challenges in comparability. Discrepancies in participant selection, validation procedures, and datasets impact the reproducibility of findings across studies \citep{ref_10_1016_j_media_2020_101694}. Many studies do not adequately address external validation within diverse populations, which may influence the generalizability of their findings \citep{ref_10_1016_j_nicl_2021_102712,ref_10_1002_alz_13390}. Overall, while significant strides have been made in leveraging AI for Alzheimer's research, the need for standardized methodologies and comprehensive validation remains critical for enhancing both the robustness and applicability of AI models.

\section{Limitations and Research Outlook}

\subsection{Paper-Reported Limitations}

\textit{This section synthesizes the reported limitations from studies applying artificial intelligence in Alzheimer's disease research, highlighting critical challenges that affect the generalizability and reproducibility of findings.}

Numerous studies within the realm of AI applications for Alzheimer's disease (AD) research exhibit similar limitations that may impact the validity and generalizability of their findings. A recurring issue centers on sample size; many investigations underscore the constraint of using small datasets, raising concerns regarding overfitting and the model's applicability across diverse populations \citep{ref_10_1016_j_media_2020_101694,ref_10_1002_alz_13390,ref_10_3389_fpubh_2022_772592,ref_10_1016_j_matpr_2021_03_061}. Additionally, methodological discrepancies, including variations in participant selection and validation procedures, complicate performance comparisons across different models \citep{ref_10_1016_j_media_2020_101694}. Certain studies noted challenges in establishing reproducible frameworks, often attributed to insufficient documentation of implementation parameters \citep{ref_10_1016_j_media_2020_101694,ref_10_3389_fpubh_2022_834032}. Generalizability is further hampered by the lack of diversity in study participants, particularly regarding demographic representation, leaving models at risk of bias \citep{ref_10_1109_access_2021_3090474,ref_10_1016_s2589_7500_22_00169_8}. Studies also highlight difficulties in employing certain neuroimaging modalities, noting the potential limitations in transferring findings from one imaging technique to another \citep{ref_10_3389_fpubh_2022_834032,ref_10_1016_j_cmpb_2021_106032}. Furthermore, several studies lacked robust validation phases and relied predominantly on cross-sectional data, which may lead to unreliable conclusions about clinical outcomes \citep{ref_10_1016_j_neuroimage_2020_117203,ref_10_1002_alz_12032} and raised questions regarding potential biases in participant selection \citep{ref_10_1016_j_nicl_2021_102712}.

\subsection{Explicit Research Gaps}

\textit{This section delves into specific research gaps identified in recent studies concerning the application of artificial intelligence in Alzheimer's disease research. It highlights the critical need for developing standardized frameworks, improving data diversity, and enhancing model validation across varied populations.}

Recent investigations into artificial intelligence applications in Alzheimer's disease have identified several explicit research gaps warranting further exploration. Notably, there is a pressing need for standardized frameworks to enhance reproducibility in convolutional neural network (CNN) research \citep{ref_10_1016_j_media_2020_101694}. Future work is encouraged in applying AI technologies across diverse imaging modalities and larger datasets to improve generalizability in diagnostic models \citep{ref_10_3389_fpubh_2022_834032}. Moreover, challenges surrounding dataset suitability underscore the necessity for AI approaches to adequately address complexities in data interaction \citep{ref_10_1002_alz_13390}. Further studies suggest incorporating additional biomarkers into AI models to bolster detection accuracy \citep{ref_10_1016_s2589_7500_22_00169_8}. Furthermore, validating models across diverse population cohorts remains crucial to ensure their comprehensive applicability and effectiveness \citep{ref_10_1002_alz_12032,ref_10_1016_j_nicl_2021_102712}. These identified gaps highlight the urgent requirement for continued research focused on enhancing model performance and establishing data-driven frameworks that incorporate a variety of biological and clinical factors, ultimately leading to improved predictive and diagnostic tools for Alzheimer's disease.

\subsection{Future Research Directions}

\textit{This section synthesizes future research directions for the application of artificial intelligence (AI) in Alzheimer's disease, as explicitly stated by recent studies. It identifies key areas underscored by the literature, including the need for standardized frameworks, diverse datasets, and novel techniques in AI methodologies to enhance early diagnosis and treatment of Alzheimer's disease. Furthermore, challenges such as data validation across populations and integration of multimodal data are highlighted as essential for improving research outcomes in this domain.}

Recent studies indicate several imperatives for future research concerning the application of artificial intelligence (AI) in Alzheimer's disease. A significant focus should be on establishing standardized frameworks to enhance reproducibility in studies involving convolutional neural networks (CNNs), as noted by \citet{ref_10_1016_j_media_2020_101694}. There is a pressing need to explore applications of AI across various imaging modalities and work with larger, more diverse datasets to improve model generalizability \citep{ref_10_3389_fpubh_2022_853294,ref_10_3389_fpubh_2022_834032}. Moreover, advancing AI techniques must address the complexities associated with dataset suitability and improve data integration measures to facilitate biomarker discovery, emphasized by \citet{ref_10_1002_alz_13390}. Additionally, enhancing model performance in diverse cohorts and integrating additional biomarkers are pivotal avenues for future exploration \citep{ref_10_1016_j_nicl_2021_102712,ref_10_1016_s2589_7500_22_00169_8}. Incorporating fuzzy logic and multimodal data fusion could not only improve diagnostic accuracy but also deepen our understanding of disease progression \citep{ref_10_1109_tfuzz_2024_3409412,ref_10_1109_icscds53736_2022_9760858}. A collaborative approach that consolidates these dimensions will be crucial for the evolution of effective AI applications in Alzheimer's research.

\section{Conclusion}

\textit{This literature review synthesizes significant advancements in applying artificial intelligence (AI) to Alzheimer's disease research, highlighting the strengths of various methodologies while acknowledging existing limitations and knowledge gaps.}

This literature review integrates findings from 83 relevant articles published since 2020, demonstrating that AI, particularly through deep learning and machine learning techniques, plays a critical role in understanding, diagnosing, and treating Alzheimer's disease. Notable advancements include the use of convolutional neural networks (CNNs) to classify Alzheimer’s patients with high accuracy \citep{ref_10_1016_j_media_2020_101694,ref_10_3389_fpubh_2022_853294}, as well as models that effectively predict disease progression \citep{ref_10_1016_j_neuroimage_2020_117203,ref_10_3389_fpubh_2022_834032}. Despite these advancements, several challenges remain, particularly regarding model generalizability across diverse populations and datasets \citep{ref_10_1002_alz_12032,ref_10_1002_alz_13390}. Limitations also arise from variability in methodologies employed across studies, complicating comparative assessments and hindering reproducibility \citep{ref_10_1016_j_media_2020_101694}. Gaps in research include the need for comprehensive validation of models against broader datasets and integration of multimodal data sources \citep{ref_10_1016_j_cmpb_2021_106032,ref_10_1007_s44174_023_00078_9}. Overall, while AI offers promising tools for Alzheimer's disease research, future research must address methodological inconsistencies and strive toward standardized frameworks for validation and reproducibility in the field.

\clearpage
\bibliographystyle{abbrvnat}
\nocite{ref_10_1016_j_patcog_2022_109031,ref_10_3389_fpubh_2022_834032,ref_10_1016_j_compbiomed_2022_106240,ref_10_1016_j_jbi_2020_103411,ref_10_1016_s2589_7500_22_00169_8,ref_10_1002_alz_12032,ref_10_1016_j_neurobiolaging_2021_04_024,ref_10_1093_jamia_ocac168,ref_10_1016_j_nicl_2021_102712,ref_10_1016_j_irbm_2020_06_006,ref_10_1016_j_arr_2024_102497,ref_10_1016_j_compbiomed_2021_104678,ref_10_3389_fpubh_2022_853294,ref_10_1016_j_artmed_2022_102332,ref_10_1111_nan_12759,ref_10_1016_j_patcog_2020_107247,ref_10_1016_j_neucom_2020_01_053,ref_10_1016_j_cmpb_2021_106032,ref_10_3389_fpubh_2022_772592,ref_10_1016_j_compbiomed_2021_105032,ref_10_1016_j_compbiomed_2021_105056,ref_10_3390_brainsci10020084,ref_10_1016_j_bspc_2021_103293,ref_10_1002_alz_70175,ref_10_1109_access_2021_3090474,ref_10_1109_icaccs48705_2020_9074248,ref_10_1016_j_neuroimage_2020_117203,ref_10_1016_j_sbi_2024_102776,ref_10_1016_j_measurement_2020_108838,ref_10_1109_jbhi_2021_3053568,ref_10_1109_icscds53736_2022_9760858,ref_10_1016_j_matpr_2021_03_061,ref_10_1109_tfuzz_2024_3409412,ref_10_1007_s44174_023_00078_9,ref_10_1016_j_media_2020_101694,ref_10_1002_alz_13390,ref_10_1016_j_ejrad_2023_110934,ref_10_1016_j_mri_2021_02_001,ref_10_5582_irdr_2023_01091,ref_10_1016_j_compbiomed_2023_107328}
\bibliography{references}

\section*{Acknowledgments}
This document was generated from structured literature-review evidence produced by the AD literature knowledge pipeline.

\end{document}
